\subsection{An actual third norm level and the degree-twenty-seven class group}

\begin{statement}{A.176}{Proved}
{Support repair, the first balanced candidate, and the full third-level groups}
Retain \(K_n,E_n,S,\beta_1,\beta_2,h_{11},h_{59}\) from
A.162--A.175.  Put
\[
       b_1=\beta_1h_{11}^3h_{59}^6,\qquad
       b_2=\beta_2h_{11}h_{59}^2.
\tag{A.176.1}
\]
There exists an actual \(S_3\)-unit \(b_3\in K_3^\times\) with
\[
 N_{3,2}b_3=b_2,\qquad N_{2,1}b_2=b_1,\qquad
                         N_{1,0}b_1=-2.
\tag{A.176.2}
\]
Thus this is a norm-compatible chain of actual elements, not only
a chain of completed classes.

For the actual three-primary classes of A.174 one has
\[
       P=-3R_{59},\qquad R_{11}=5R_{59}.
\tag{A.176.3}
\]
In particular the first fractional ideal in A.174.5 has zero
three-primary class.  The complete third-level groups are
\[
 \begin{gathered}
 \operatorname{Cl}(E_3)\{3\}\simeq
                \Z/9\oplus(\Z/3)^2,\qquad
 \operatorname{Cl}(E_3)\{3\}^{\tau=+1}\simeq\Z/9,\qquad
 \operatorname{Cl}(E_3)\{3\}^{\tau=-1}\simeq(\Z/3)^2,\\
 \operatorname{Cl}(K_3)\{3\}=\langle R_{59}\rangle\simeq\Z/9,
                    \qquad C_3^{\rm bal}=0.
 \end{gathered}
\tag{A.176.4}
\]
The distinguished two-prime has exact order three in this
three-primary class group.  These conclusions concern the stated
third level.  They do not assert a compatible global norm family
through every higher level.
\end{statement}

\begin{proof}
\emph{An actual cyclic norm with its support removed.}
The fields
\[
 \begin{matrix}
 E_3=K_3(\zeta_3)&\supset&E_2=K_2(\zeta_3)\\
 \cup&&\cup\\
 K_3&\supset&K_2
 \end{matrix}
\tag{A.176.5}
\]
have vertical degree two and horizontal degree three.
The upper extension is cyclic, generated by adjoining a cube root
of \(\rho_2\), and is unramified outside \(S\).
All these are the fixed field embeddings of the radical tower.
A.162 proves that the actual element \(b_2\) is a local norm
from \(K_3/K_2\) at every place.
Base-changing each finite local algebra in that norm equation
shows that it is a local norm from \(E_3/E_2\) at every place
of \(E_2\).  This uses the determinant definition of the norm,
which commutes with scalar extension.
The cyclic Hasse norm theorem therefore supplies an actual
\(z\in E_3^\times\) with
\[
                         N_{E_3/E_2}z=b_2
\tag{A.176.6}
\]
\cite[Chapter VIII, Theorem 3.1, printed p.~234]{MilneCFT}.
The prime-to-three support correction below is what turns this
field norm into an \(S\)-unit norm.

Let \(D_z=\operatorname{div}_{S^c}(z)\).
Its ideal norm is zero because \(b_2\) is an \(S\)-unit.
For an outside-\(S\) prime the cyclic extension is unramified.
An inert prime is a one-element orbit with residue-degree-three
norm map, and hence its coefficient in \(D_z\) is zero.
On a split three-element orbit the norm is the sum of the three
integer coefficients.  Every vector with sum zero is the image
of \(\sigma-1\) on the free orbit lattice.
Choosing these preimages on the finitely many nonzero orbits gives
an actual finite fractional divisor \(D\) with
\[
                         D_z=(\sigma-1)D .
\tag{A.176.7}
\]
This is the required ideal version of Hilbert 90, with its exact
outside-\(S\) divisor map.

A.175 proves \(\operatorname{Cl}_{S}(E_3)\{3\}=0\).
Thus the class of \(D\) in the finite \(S\)-class group has order
\(h\) prime to three.  Choose \(y\in E_3^\times\) with
\(\operatorname{div}_{S^c}(y)=hD\).
Then the literal element
\[
              \epsilon=\frac{z^h}{\sigma(y)/y}
\tag{A.176.8}
\]
is an \(S\)-unit and has norm \(b_2^h\).
The norm of \(\sigma(y)/y\) is one by the cyclic product.
Choose integers \(a,b\) with \(ah+3b=1\).
The element
\[
                       z_S=\epsilon^a b_2^b
\tag{A.176.9}
\]
is an actual \(S\)-unit of \(E_3\), and its norm is exactly
\(b_2\).  Negative exponents cause no change of support.
No completion or choice of a fractional root has entered
A.176.6--9.

\emph{The quadratic norm transfer to the specified pure field.}
Put \(\gamma=N_{E_3/K_3}z_S\).
Transitivity in the actual field square A.176.5 gives
\[
 N_{K_3/K_2}\gamma
   =N_{E_2/K_2}N_{E_3/E_2}z_S
   =N_{E_2/K_2}b_2=b_2^2 .
\tag{A.176.10}
\]
Consequently
\[
                              b_3=\gamma^2/b_2
\tag{A.176.11}
\]
is an actual \(S_3\)-unit and has norm \(b_2^{4-3}=b_2\).
The remaining two norm equalities are the literal identities
of A.162.  This proves A.176.2.

\emph{The exact divisor of this norm preimage.}
The unique three-prime has residue-degree-one norm in
\(K_3/K_2\).  Since \(b_2\) has valuation zero there,
so does \(b_3\).
At each changed-prime branch the extension is totally ramified
of degree three and the ideal norm has residue weight one.
Since \(h_q\), embedded from \(K_1\), has divisor
three times the corresponding balanced \(q\)-divisor in \(K_2\),
the changed-prime part of \(\operatorname{div}(b_3)\) is exactly
\[
                         3\mathfrak R_{11}+6\mathfrak R_{59}.
\tag{A.176.12}
\]
At two, the old distinguished prime splits into the branches
of absolute residue degrees one and two; its norm equation is
\(k_0+2k_2=1\).
The old degree-two and degree-six primes become single unramified
cubic branches.  Their target valuations are zero, so their
upper valuations are zero.
For some integer \(k\), therefore,
\[
 \operatorname{div}(b_3)=
 (1-2k)P_{2,0,3}+kP_{2,2,3}
                 +3\mathfrak R_{11}+6\mathfrak R_{59}.
\tag{A.176.13}
\]
All other valuations vanish by the proved \(S\)-unit property.
The actual principal divisor of the embedded \(\beta_2\) is
\(P_{2,0,3}+P_{2,2,3}\), by A.174.16.
Also \(3(R_{11}+R_{59})=0\), with the literal principal
function of A.174.14.
Passing A.176.13 to the three-primary class group thus gives
\[
                         (1-3k)P+3R_{59}=0 .
\tag{A.176.14}
\]
The integer \(1-3k\) is a three-adic unit, and \(3R_{59}\)
has order three by A.174.
Its inverse acts as one on this order-three subgroup.
It follows that \(P=-3R_{59}\).
This proves the first equality of A.176.3 and the vanishing
of the first specified candidate; it also proves that \(P\)
has exact order three.

\emph{All third-level ordinary classes and their eigenspaces.}
Write \(A_3=\operatorname{Cl}(E_3)\{3\}\).
A.175 gives
\[
                   |A_3^{\langle\sigma\rangle}|=81,
                  \qquad\operatorname{Cl}_{S}(E_3)\{3\}=0.
\tag{A.176.15}
\]
The latter assertion means, through the ordinary ideal class
exact sequence, that \(A_3\) is generated by the three-primary
classes of the primes over \(S\).
We check the action on every such generator.
The three-prime and all six changed-prime branches are totally
ramified in \(E_3/E_2\), so their prime ideals are fixed by
\(\sigma\).
Four of the five two-primes of \(E_2\) are inert unramified
cubic primes in \(E_3\), and their upper primes are also fixed.
The remaining old two-prime splits into exactly three primes,
permuted by \(\sigma\).

One member \(\mathfrak p_0\) of this last orbit lies above the
distinguished degree-one prime \(P_{2,0,3}\) of \(K_3\).
The quadratic extension \(E_3/K_3\) is unramified of residue
degree two there, so ideal extension sends
\(P_{2,0,3}\) to precisely \(\mathfrak p_0\), with ramification
weight one.
The established equality \(P=-3R_{59}\) therefore expresses
\([\mathfrak p_0]_{(3)}\) as a combination of changed-prime
classes.  All those classes are fixed by \(\sigma\).
Thus the class of \(\mathfrak p_0\) is fixed, and its other
two conjugate prime classes are the same class.
Every generator of \(A_3\) is now fixed, so
\[
                              A_3=A_3^{\langle\sigma\rangle},
                              \qquad |A_3|=81 .
\tag{A.176.16}
\]

The actual ideal norm \(A_3\to A_2=\operatorname{Cl}(E_2)\{3\}\)
is surjective.  By A.173, \(A_2\simeq(\Z/3)^3\) is generated
by changed-prime classes; every one is the norm of its unique
totally ramified upper prime, whose residue norm weight is one.
Therefore the three-rank of \(A_3\) is at least three.
Quadratic ideal extension from \(K_3\) to \(E_3\) is injective
on three-primary classes, because its composite with norm is
multiplication by two.  Hence A.174's order-nine class \(R_{59}\)
gives an order-nine element in \(A_3\).
The classification of finite abelian groups of order \(3^4\)
now gives
\[
                             A_3\simeq\Z/9\oplus(\Z/3)^2.
\tag{A.176.17}
\]

The commuting ideal norm and involution maps send the minus
eigenspace of \(A_3\) onto the minus eigenspace of \(A_2\),
which is \((\Z/3)^2\) by A.173.
For example, project any norm preimage with the actual
idempotent \((1-\tau)/2\); division by two is unique on
these finite three-groups.
Thus \(|A_3^-|\ge9\).
The plus eigenspace contains the extended \(R_{59}\) of
order nine, so \(|A_3^+|\ge9\).
Their product has order \(81\), forcing
\[
                              A_3^+\simeq\Z/9,\qquad
                              A_3^-\simeq(\Z/3)^2.
\tag{A.176.18}
\]
For the actual quadratic extension the ideal extension \(j\)
and norm \(N\) satisfy
\[
                   Nj=2,\qquad jN=1+\tau.
\tag{A.176.19}
\]
The second identity includes the ramification multiplicities:
at a ramified quadratic prime, extension after norm is twice
that fixed prime, exactly as on the right.
Consequently \(j\) identifies
\(\operatorname{Cl}(K_3)\{3\}\) with \(A_3^+\).
The class \(R_{59}\) generates this cyclic order-nine group.
The character \(\theta\) of A.174 is therefore an isomorphism
on it.  Its specified value on \(R_{11}\) gives
\(R_{11}=5R_{59}\), proving the remaining equality in A.176.3.
The balanced quotient kills \(R_{59}\) itself, so it is zero.

All support and coefficient maps used here are finite-level maps.
The actual element \(b_3\) proves the displayed norm chain, and
the complete class calculation proves the entire balanced
third-level test.  No transition to level four or to the global
inverse limit is inferred from these conclusions.
The finite orbit-divisor, norm-exponent and class-group
constraints are checked in
\path{code/two_inert_level3_support_repair_certificate.py}.
\end{proof}
